\section{The actual product local frame and the original Gauss source frame}

\begin{definition}\label{def:fc-point-laurent-local-open-factor}
\lean{EnvelopeProperty.Berkovich.SchemePointLaurentLocalFrame.completedOpenFactor,EnvelopeProperty.Berkovich.SchemePointLaurentLocalFrame.completedOpenFactor_comp,EnvelopeProperty.Berkovich.SchemePointLaurentLocalFrame.productOpenFactor,EnvelopeProperty.Berkovich.SchemePointLaurentLocalFrame.productOpenFactor_comp,EnvelopeProperty.Berkovich.SchemePointLaurentLocalFrame.productOpenFactor_restrict}
\leanok
\uses{def:fc-point-completed-line-frame,cor:actual-line-model-norm-local-frame,def:native-point-laurent-gauss-family}
Let $p:X\to\operatorname{Spec}R$ be proper over a bounded valued ring,
let $x$ be an original relative valued point, and let $U$ be an
original open neighborhood of its true proper center. The actual
completed point factors through $U$ by composing the original
residue embedding with the original proper-center generic factor.
Let $g:X_S\to X$ and $b:\operatorname{Spec}H(x)[T,T^{-1}]\to X_S$
be the actual original Laurent projection and Gauss product map.
Then $b$ factors through the actual original open $g^{-1}U$, and
its restricted projection to $U$ is the constant-field projection
followed by the completed-point factor just constructed.
\end{definition}
\begin{proof}
The original residue generic triangle gives the completed-point
triangle after composition with the residue embedding. The actual
product first-projection triangle then shows that every original
source point maps into $g^{-1}U$. The open-immersion lifting property
constructs the whole factor of $b$ into this original open. Its
composite with the original open inclusion is exactly $b$. To
compare the two maps to $U$, compose them with the original open
inclusion $U\to X$. The product first-projection triangle and the
completed-point triangle make these composites equal. The open
inclusion is a monomorphism, proving the asserted whole restricted
projection identity.
\end{proof}

\begin{theorem}\label{thm:fc-point-laurent-native-local-frame}
\lean{EnvelopeProperty.Berkovich.SchemePointLaurentLocalFrame.completedFrame_eq,EnvelopeProperty.Berkovich.SchemePointLaurentLocalFrame.sourceFrame_eq_localPullbackFrame}
\leanok
\uses{def:fc-point-laurent-local-open-factor,def:fc-point-completed-line-frame,def:fc-point-laurent-line-evaluation,lem:actual-line-frame-via-factor-composition,lem:fc-original-section-pullback-frame,alggeom:actual-original-frame-pullback-composition}
An actual original frame of $L$ on $U$ pulls back to its actual
original product frame on $g^{-1}U$. Pull this product frame through
the whole source factor just constructed. The resulting entire
original source frame is exactly the source frame constructed from
the same proper-center completed-point frame in the actual
Laurent section-evaluation theorem. No frame-preservation or
coefficient-matching identity is supplied.
\end{theorem}
\begin{proof}
The native frame-composition identity first identifies the completed
point frame with direct transport of the original $U$ frame through
the completed-point factor. Applying the same composition identity
through the constant-field projection identifies its source frame
with direct transport through the composite source-to-$U$ map,
preceded by the native product pullback-composition comparison.
For the other frame, native compatibility of frame transport with
pullback composition moves the product comparison past the source
open factor. The frame-composition identity then transports its
original $U$ frame through the restricted source-to-$U$ composite.
The preceding whole restricted projection equality identifies these
two actual composites. Thus the whole original frame isomorphisms
are equal, including all native composition and unit comparisons.
\end{proof}

\begin{theorem}\label{thm:fc-point-laurent-local-coefficient-value}
\lean{EnvelopeProperty.Berkovich.SchemePointLaurentLocalFrame.localCoefficient,EnvelopeProperty.Berkovich.SchemePointLaurentLocalFrame.sourceFrameSection_eq_localCoefficient,EnvelopeProperty.Berkovich.SchemePointLaurentLocalFrame.openPoint,EnvelopeProperty.Berkovich.SchemePointLaurentLocalFrame.openPoint_comp,EnvelopeProperty.Berkovich.SchemePointLaurentLocalFrame.localCoefficient_value,EnvelopeProperty.Berkovich.SchemePointLaurentLocalFrame.localCoefficient_finite_family}
\leanok
\uses{thm:fc-point-laurent-native-local-frame,thm:fc-point-laurent-line-gauss,thm:fc-point-completed-frame-model-norm,cor:fc-point-exists-completed-frame-model-norm,scheme:valued-points,lem:fc-original-line-section-fibre}
Let $L$ be an original invertible line and fix its actual original
proper-center frame on $U$. The actual native product Gauss point
has a genuine native lift to $g^{-1}U$. For every original whole
section $\sigma$ of $g^*L$, pull it back to this original open and
apply the actual original product frame there. Evaluate this true
local coefficient at that native lift. Its value is exactly
\[
 \max_i \|m_i\|_{L,x}r^i,
\]
where $m_i$ are the unique actual original Laurent coefficients
and $\|\cdot\|_{L,x}$ is the true original model line norm of the
original native residue pullback-unit section. For every finite
original section family, the same actual original finite coefficient
and integer-weight family gives its maximum before all points,
frames and positive radii. All negative exponents and zero values
are included.
\end{theorem}
\begin{proof}
Map the actual source Gauss point through the constructed open
factor. Its composite through the original product-open inclusion
is the original product Gauss point, by the whole factor triangle
and genuine native valued-point functoriality. The native
pullback-unit section is the value at one of its genuine pulled
section morphism. Applying an actual frame to this whole morphism
gives its actual scalar endomorphism. Pulling this scalar through
the open factor, native section naturality and the preceding whole
frame identity show that the true product local coefficient pulls
back to exactly the actual source-frame section. Native valued-point
evaluation naturality identifies their actual values. The original
source Gauss formula computes this value as the finite weighted
maximum of the completed-frame functional values. Their absolute
values equal the true original model line norms by the proper-center
completed-frame theorem. Substitution gives the asserted formula.
The finite-family source formula and the same substitution give
the common original pair family, whose construction has no point,
frame, valuation or radius input.
\end{proof}
